\section{Helsetjenestens driftsorganisasjon for nødnett HF}
\label{sec:hdo}
Helsetjenestens driftsorganisasjon for nødnett HF (HDO) ble 
stiftet 29. april 2013 og har hovedsete på Gjøvik \cite{hdo_om_oss}. Foretaket 
eies av de fire regionale helseforetakene og er organisert 
etter helseforetaksloven \cite{hdo_eierstyring}.

HDOs samfunnsoppdrag er å levere kommunikasjons- og 
beslutningsstøtteløsninger til den akuttmedisinske kjede, 
inkludert drift av nødnummer 113, legevaktnummer 116~117 
og løsninger for video- og datadeling \cite{hdo_om_oss}. 
HDO har eier-, drifts- og forvaltningsansvar for alt 
brukerrelatert utstyr tilknyttet nødnett i helse, og drifter 
166 kontrollrom, herunder AMK-sentraler, legevaktsentraler 
og akuttmottak, over hele Norge \cite{hdo_kontrollrom}.

En av disse tjenestene er videoløsningen. Når en innringer 
kontakter en AMK- eller legevaktsentral per telefon, kan 
operatøren sende en SMS-lenke som gir innringeren mulighet 
til å starte en enveis videostrøm tilbake til sentralen. Lyd 
går via den ordinære telefonsamtalen; video brukes som 
beslutningsstøtte for å gi operatøren bedre 
situasjonsforståelse. Løsningen lagrer ingen lyd, bilde 
eller pasientopplysninger \cite{hdo_produktark}.

HDOs oppdragsdokument for 2026 viderefører ansvaret for 
videokommunikasjon i akuttmedisinsk kjede og pålegger HDO 
å bidra til interoperabilitet og økt standardisering av 
videoløsninger \cite{hdo_oppdrag2026}. VideoAPI-prosjektet 
ble initiert som et konkret bidrag til dette arbeidet.